\chapter{Projective Varieties}

% -------------------- % -------------------------------------------
% 4.1 PROJECTIVE SPACE % -------------------------------------------
% -------------------- % -------------------------------------------
\section{Projective Space}

\begin{definition}
	Let $k$ be any field. We define the \textit{projective $n$ space} over $k$, written as $\mbb{P}^{n}(k)$ (or simply $\mbb{P}^{n}$), with the set of equivalence classes of points in $\mbb{A}^{n+1} \setminus \{(0,\dots,0)\}$, where two points $(x),(y) \in \mbb{A}^{n+1} \setminus \{(0,\dots,0)\}$ are $(x) \sim (y)$ if and only if there is a nonzero $\lambda \in k$ such that $x_{i} = \lambda y_{i}$ for all $i = 1,\dots,n+1$. 
\end{definition}

\begin{remark}
	It's easy to see that we are relating points if they are on the same (unique) line that passes through the origin on $\mbb{A}^{n+1}$. Then, points in $\mbb{P}^{n}$ are lines passing through the origin in $\mbb{A}^{n+1}$.
\end{remark}

\begin{definition}
	We let $U_{i} = \{[x_{1} : \ldots : x_{n+1}] \in \mbb{P}^{n} \ | \ x_{i} \neq 0\}$. Each $P \in U_{i}$ can be written uniquely in the form
	\[ P = [ x_{1} : \ldots : x_{i} : \ldots : x_{n+1} ]. \]
	Then, we define $\varphi_{i} \colon \mbb{A}^{n} \to U_{i}$ by $\varphi_{i}(a_{1},\ldots,a_{n}) = [a_{1} : \ldots : a_{i} : \ldots : a_{n}]$, which defines a one-to-one correspondence between the points of $\mbb{A}^{n}$ and the points of $U_{i}$.
\end{definition}

\begin{remark}
	Note that $\mbb{P}^{n} = \bigcup_{i=1}^{n} U_{i}$, so $\mbb{P}^{n}$ is covered by $n+1$ sets each of which look just like affine $n$-space.
\end{remark}

\begin{definition}
	Let
	\[ H_{\infty} \coloneq \mbb{P}^{n} \setminus U_{n+1} = \{[x_{1} : \ldots : x_{n+1}] \ | \ x_{n+1} = 0\}; \]
	$H_{\infty}$ is often called the \textit{hyperplane at infinity}.
\end{definition}

\begin{remark}
	The correspondence $[x_{1} : \ldots : x_{n} : 0] \leftrightarrow [x_{1} : \ldots : x_{n}]$ shows that $H_{\infty}$ may be identified with $\mbb{P}^{n+1}$. Thus, $\mbb{P}^{n} = U_{n+1} \cup H_{\infty}$ is the union of an affine $n$-space and a set that gives all directions in affine $n$-space.
\end{remark}


% ----------------------------- % ----------------------------------
% 4.2 PROJECTIVE ALGEBRAIC SETS % ----------------------------------
% ----------------------------- % ----------------------------------
\section{Projective Algebraic Sets}
\begin{definition}
	A point $P \in \mbb{P}^{n}$ is said to be a \textit{zero} of a polynomial $F \in k[X_{1},\ldots,X_{n+1}]$  if 
	\[ F(x_{1},\ldots,x_{n+1}) = 0 \]
	for every choice of homogeneous coordinates $(x_{1},\ldots,x_{n+1})$ for $P$; we then write $F(P) = 0$.
\end{definition}

\begin{definition}
	For any set $S$ of polynomials in $k[X_{1},\ldots,X_{n+1}]$, we let 
	\[ V(S) = \{P \in \mbb{P}^{n} \ | \ P \text{ is a zero of each } F \in S\}. \]
\end{definition}

\begin{remark}
	If $I$ is the ideal generated by $S$, $V(I) = V(S)$.
\end{remark}

\begin{definition}
	If $I = (F^{(1)}, \ldots, F^{(r)})$, where $F^{(i)} = \sum F_{j}^{(i)}$, $F_{j}^{(i)}$ a form of degree $j$, then $V(I) = V(\{F_{j}^{(i)}\})$, so $V(S) = V(\{F_{j}^{(i)}\})$ is the set of zeros of a finite number of forms. Such a set is called an algebraic set in $\mbb{P}^{n}$, or a \textit{projective algebraic space}.
\end{definition}

\begin{definition}
	For any set $X \subset \mbb{P}^{n}$, we let 
	\[ I(X) = \{F \in k[X_{1},\ldots,X_{n+1}] \ | \ \text{every $P \in X$ is a zero of $F$}\}. \]
	The ideal $I(X)$ is called the \textit{ideal} of $X$. 
\end{definition}

\begin{definition}
	An ideal $I \subset k[X_{1},\ldots,X_{n+1}]$ is called \textit{homogeneous} if for every $F = \sum_{i=1}^{m} F_{i} \in I$, $F_{i}$ a form of degree $i$, we have also $F_{i} \in I$. 
\end{definition}

\begin{remark}
	For any set $X \subset \mbb{P}^{n}$, $I(X)$ is a homogeneous ideal.
\end{remark}

\begin{prop}
	An ideal $I \subset k[X_{1},\ldots,X_{n+1}]$ is homogeneous if and only if it is generated by a (finite) set of forms.
\end{prop}

% ----------------------------------- % ----------------------------
% 4.3 AFFINE AND PROJECTIVE VARIETIES % ----------------------------
% ----------------------------------- % ----------------------------
\section{Affine and Projective Varieties}



% ------------------------- % --------------------------------------
% 4.4 MULTIPROJECTIVE SPACE % --------------------------------------
% ------------------------- % --------------------------------------
\section{Multiprojective Space}
